\section{Selecting coefficient tests and omitted sets}
\label{l:domains-section}

The common coset truncation of Section~\ref{l:core-section} already
applies to the marginal hypotheses considered here. This section gives
two different constructions of a retained measure. With a fixed
partition, we delete points detected by large coefficient tests and
control the union of all those tests. With an average over omitted sets,
we instead retain permutations for which most restrictions have bounded
density. The good omitted sets then depend on the permutation; we do not
select a fixed partition on which the retained law has all its coset
caps. Recovering a vector from their subgroup projections is the new
step, because projections for overlapping sets need not commute.

Write $[n]=\{1,\ldots,n\}$ and let $U_n$ be uniform on $S_n$.
For $R\subseteq[n]$, put $H_R=\Sym(R)$, acting trivially outside $R$,
and write $g_{-R}=g|_{[n]\setminus R}$. Its fibers are the left cosets
$gH_R$: two permutations have the same restriction precisely when
they differ by right composition with an element of $H_R$.
For a probability $\mu$ on $S_n$, let $\mu_{-R}$ and $U_{-R}$ be
the restriction laws under $\mu$ and $U_n$, respectively.
Every average over $R$ below is uniform over subsets of its stated size.

\subsection{Pruning coefficient tests for a fixed partition}

The first argument uses each marginal estimate once, on a union of
tests chosen during a finite deletion procedure. Applying the marginal
estimate separately at each deletion would multiply its error by the
number of tests and lose the conclusion.

\begin{theorem}[Spectral pruning]
\label{l:spectral-pruning}
Let $n\to\infty$ through multiples of $16$. For each $n$, partition
$[n]$ into $b=16$ equal blocks $R_i$, and let $\mu_n$ be a probability
on $S_n$ satisfying
\[
 \|(\mu_n)_{-R_i}-U_{-R_i}\|_{\TV}\le\delta_n
 \quad(1\le i\le16),\qquad \delta_n\longrightarrow0.
\]
There is a nonnegative measure $0\le\nu_n\le\mu_n$ of mass
$1-\gamma_n$, with $\gamma_n\to0$, such that
\[
 \|\widehat\nu_n(\rho)\|_{\op}\le D^{-1/8}
 \quad\text{for every irreducible $\rho$ of degree $D>1$}.
\]
The operator bound has constant one, and $\nu_n$ is not normalized.
\end{theorem}

\begin{proof}
Put $a=1/8$ and $K_i=H_{R_i}$. Fix an irreducible $\rho$ of degree
$D>1$. Every tensor type in its restriction to $\prod_iK_i$ has
product of factor degrees at most $D$. Assign each entire isotypic
space to one factor whose degree is at most $D^{1/b}$. This yields
an orthogonal decomposition of any unit vector $v$ as
$v=\sum_i v_i$, where $v_i$ contains only these small $K_i$-types.
For $W_i=\operatorname{span}\rho(K_i)v_i$,
Lemma~\ref{l:single-orbit} counts each type once, including all its
multiplicities. Together with Lemma~\ref{l:degree-inverse-first}, it gives
\[
 \dim W_i\le C_1D^{3/b},\qquad
 C_1=\sup_{m\ge1}\sum_{\tau\in\widehat{S_m}}(\dim\tau)^{-1}.
\]
For fixed unit $u,v$, the score
\[
 H_{\rho,u,v,i}(g)=\|\operatorname{Proj}_{\rho(g)W_i}u\|
\]
depends only on $g_{-R_i}$. Its uniform second moment is
$\dim W_i/D$ by Schur averaging
\cite[Proposition~3.14]{etingof}. Consequently
\begin{equation}
 U_n\{H_{\rho,u,v,i}>D^{-a}/(2b)\}
 \le4b^2C_1D^{-1+2a+3/b}.
 \label{l:pruning-test-bound}
\end{equation}
If $|\langle u,\rho(g)v\rangle|>D^{-a}/2$, at least one score
exceeds $D^{-a}/(2b)$, because $\|v_i\|\le1$ and
$|\langle u,\rho(g)v_i\rangle|\le H_{\rho,u,v,i}(g)$.

Start with the measure $\nu=\mu_n$. Whenever its claimed operator
bound fails, choose $\rho,u,v$ with
$|\langle u,\widehat\nu(\rho)v\rangle|>D^{-a}$, record this
test, and delete all remaining mass on
\[
 \{g:|\langle u,\rho(g)v\rangle|>D^{-a}/2\}.
\]
The mass deleted is greater than $D^{-a}/2$: outside this event the
coefficient has absolute value at most $D^{-a}/2$, everywhere it has
absolute value at most one, and the current measure has mass at most
one. Since all deleted pieces are disjoint, any one representation of
degree $D$ can be selected at most $2D^a$ times. There are finitely many
irreducibles, so the procedure terminates with all the required bounds.
The choices are made from deterministic current measures; thus the
eventual finite list of tests is fixed, not selected anew for a sampled
permutation.

For each $i$, let $E_i$ be the union of its score events in
\eqref{l:pruning-test-bound} over the complete list of recorded tests.
It is a single event determined by $g_{-R_i}$. Counting tests for each
representation gives
\[
 U_n(E_i)\le8b^2C_1\sum_{\rho:\dim\rho>1}
             D^{-1+3a+3/b}
 =8b^2C_1\sum_{\rho:\dim\rho>1}D^{-7/16}=o(1),
\]
by Theorem~\ref{l:degree-all-powers}. Apply the marginal hypothesis to
this entire union to get $\mu_n(E_i)\le U_n(E_i)+\delta_n$.
Every deleted point belongs to some $E_i$, so the total mass loss is
at most $\sum_i\mu_n(E_i)=o(1)$. This proves the theorem.
\end{proof}

\subsection{A common truncation and its low-dimensional types}

We first give the truncation and the projection estimates with a variable
number of omitted blocks. This identifies exactly which part of the two
transfers below is shared.

\begin{lemma}[Good domains and averaged projections]
\label{l:adaptive-domain-preparation}
Let $q\ge2$ divide $n$, put $r=n/q$, and let $\mu$ be a probability
on $S_n$. Define
\[
 \delta=\E_{|R|=r}\|\mu_{-R}-U_{-R}\|_{\TV},\qquad
 e_R(g)=\mathbf1\{\mu(gH_R)\le2/[S_n:H_R]\},
\]
and
\[
 E=\{g:\E_R e_R(g)\ge7/8\}.
\]
Then $\mu(E)\ge1-16\delta$. If $\mu(E)\ge1/2$, the probability
$\nu=\mu(\,\cdot\mid E)$ satisfies
$\|\nu-\mu\|_{\TV}=1-\mu(E)\le16\delta$.

Fix an irreducible unitary representation $\rho$ of $S_n$ on $V$,
with $D=\dim V>1$, and fix $\xi>0$. Let $P_R$ project onto the sum
of the $H_R$-isotypic spaces whose irreducible type has degree at most
$D^\xi$. Then
\begin{equation}
 \E_R P_R=cI,\qquad c\ge1-\frac1{q\xi}.
 \label{l:adaptive-scalar-projection}
\end{equation}
For $z\in V$ with $\|z\|\le1$, put
$W_R(z)=\operatorname{span}\{\rho(h)P_Rz:h\in H_R\}$.
With the absolute constant
\[
 C_2=\sup_{m\ge1}\sum_{\tau\in\widehat{S_m}}(\dim\tau)^{-2}
\]
from Lemma~\ref{l:degree-inverse-square}, one has
\begin{equation}
 \dim W_R(z)\le C_2D^{4\xi}.
 \label{l:adaptive-small-orbit}
\end{equation}
If $\mu(E)\ge1/2$, then for every unit $u\in V$,
\begin{equation}
 \E_{g\sim\nu} e_R(g)
       |\langle u,\rho(g)P_Rz\rangle|
 \le 2\sqrt{C_2}\,D^{-1/2+2\xi}.
 \label{l:adaptive-good-coefficient}
\end{equation}
\end{lemma}

\begin{proof}
For any probability $a$ and uniform probability $b$ on a finite set,
\[
 \sum_{x:a(x)>2b(x)}a(x)
 \le2\sum_{x:a(x)>b(x)}(a(x)-b(x))
 =2\|a-b\|_{\TV}.
\]
Apply this to each restriction law. Under $g\sim\mu$, the expected
fraction of sets with $e_R(g)=0$ is at most $2\delta$. Markov's
inequality gives $\mu(E^c)\le16\delta$. Conditioning on $E$ changes
the law in total variation by exactly $\mu(E^c)$.

Conjugating $H_R$ by $a\in S_n$ replaces it by $H_{aR}$ and conjugates
$P_R$ by $\rho(a)$. Hence the average projection commutes with
$\rho(S_n)$ and is scalar. To estimate its trace, choose one partition
$(R_1,\ldots,R_q)$ into equal blocks. The restriction of $V$ to
$\prod_iH_{R_i}$ is an orthogonal sum of tensor types, with arbitrary
multiplicities. If a type has factor degrees $d_1,\ldots,d_q$, then
$\prod_i d_i\le D$, since at least one copy of that type occurs in $V$.
Consequently at most $1/\xi$ factors have degree greater than $D^\xi$.
On this tensor type, $I-P_{R_i}$ is either the identity or zero according
to that inequality. Summing traces, including multiplicities, gives
\[
 \sum_{i=1}^q\operatorname{rank}(I-P_{R_i})\le D/\xi.
\]
All these ranks agree by conjugacy. This proves
\eqref{l:adaptive-scalar-projection}. Commutation was used only for
the disjoint subgroups in this rank calculation.

For the orbit bound, Lemma~\ref{l:single-orbit} gives, with all
multiplicities included,
\[
 \dim W_R(z)\le\sum_{\dim\tau\le D^\xi}(\dim\tau)^2
 \le D^{4\xi}\sum_\tau(\dim\tau)^{-2},
\]
which is \eqref{l:adaptive-small-orbit}.

It remains to turn the orbit estimate into a coefficient estimate.
The space $W=W_R(z)$ is $H_R$-invariant, so $\rho(g)W$ depends only
on $gH_R$. For unit $u$, Schur averaging and the trace give
\[
 \E_{g\sim U_n}\|\operatorname{Proj}_{\rho(g)W}u\|^2
 =\frac{\dim W}{D}.
\]
On a coset with $e_R=1$, the $\nu$-mass is at most four times its
uniform mass: the original cap is two and normalization costs at most
two. Since
$|\langle u,\rho(g)P_Rz\rangle|\le
\|\operatorname{Proj}_{\rho(g)W}u\|$, Cauchy--Schwarz gives
\[
 \E_\nu e_R(g)|\langle u,\rho(g)P_Rz\rangle|
 \le\left(4\frac{\dim W}{D}\right)^{1/2}.
\]
This proves \eqref{l:adaptive-good-coefficient}.
\end{proof}

\subsection{A normalized expansion for 1024 omitted blocks}

Both transfers below use the event $E$ and conditional law $\nu$
constructed in the preceding lemma, with their respective values of $q$.
The first uses the small lost rank in
\eqref{l:adaptive-scalar-projection} to normalize the average of the
retained projections. Its complement is then the average of only the
discarded projections. This gives a rapidly decreasing remainder.

\begin{theorem}[Omitted-set transfer]
\label{l:omitted-set-transfer}
Let $n\to\infty$ through multiples of $1024$, and let $\mu_n$ be
probabilities on $S_n$ such that
\[
 \delta_n=\E_{|R|=n/1024}
       \|(\mu_n)_{-R}-U_{-R}\|_{\TV}\longrightarrow0.
\]
For all sufficiently large $n$, condition $\mu_n$ on the event $E$
of Lemma~\ref{l:adaptive-domain-preparation}, obtaining $\nu_n$.
Then $\|\nu_n-\mu_n\|_{\TV}\le16\delta_n$, and there is an absolute
$C$ such that every irreducible representation $\rho$ of degree $D>1$
satisfies
\[
 \|\widehat\nu_n(\rho)\|_{\op}\le CD^{-1/4}.
\]
If in addition $\E_{\mu_n}\sgn(g)=o(1)$, then
$\|\mu_n^{*8}-U_n\|_{\TV}\to0$.
\end{theorem}

\begin{proof}
Suppress $n$, fix $\rho$, and use Lemma~\ref{l:adaptive-domain-preparation}
with $q=1024$ and $\xi=1/16$. Thus $c\ge63/64$, and the
right side of \eqref{l:adaptive-good-coefficient} is
$A D^{-3/8}$, where $A=2\sqrt{C_2}$. For $g\in E$ define
\[
 Q_g=\E_R e_R(g)P_R,\qquad
 F_g=I-Q_g/c=c^{-1}\E_R(1-e_R(g))P_R.
\]
The equality for $F_g$ uses the scalar identity $\E_RP_R=cI$;
it does not require the individual projections to commute.
These are positive operators, and
$\|F_g\|_{\op}\le1/(8c)$. Put $s=\lceil\log D\rceil$, with
natural logarithms. The finite identity
\begin{equation}
 I=\sum_{j=0}^{s-1}(Q_g/c)F_g^j+F_g^s
 \label{l:omitted-finite-expansion}
\end{equation}
holds because $Q_g/c=I-F_g$.

We now bound the expected coefficient of each term while preserving its
projection order. For fixed unit vectors $u,w$, the $j$th summand in
\eqref{l:omitted-finite-expansion} expands as an average over independently
chosen $R_0,\ldots,R_j$, with scalar factor $c^{-(j+1)}$, indicator
\[
 e_{R_0}(g)\prod_{i=1}^j(1-e_{R_i}(g)),
\]
and vector $P_{R_0}P_{R_1}\cdots P_{R_j}w$. Fix the sets before
integrating over $g$. The vector $z=P_{R_1}\cdots P_{R_j}w$ is then
deterministic with norm at most one. Take absolute values and discard
the failure indicators for $i\ge1$. Equation~\eqref{l:adaptive-good-coefficient}
now bounds the expected coefficient by
$c^{-(j+1)}AD^{-3/8}$. In particular the argument never exchanges two
projections.

The remainder has norm at most $(8c)^{-s}$. Summing gives
\[
 \E_\nu|\langle u,\rho(g)w\rangle|
 \le As c^{-s}D^{-3/8}+(8c)^{-s}
 \le CD^{-1/4}.
\]
For the last inequality, $c^{-s}\le c^{-1}D^{\log(1/c)}$ and
$\log(1/c)\le\log(64/63)<1/32$; the remaining logarithmic factor is
absorbed by $D^{3/32}$. Also $\log(8c)>1/4$. Taking the supremum
over the unit vectors proves the operator bound.

To obtain the convolution conclusion, the sign coefficient of $\nu_n$
is $o(1)$, since the original sign coefficient is $o(1)$ and
$\|\nu_n-\mu_n\|_{\TV}=o(1)$. Plancherel and
$\|B\|_{\HS}\le\sqrt D\|B\|_{\op}$ yield
\[
 4\|\nu_n^{*8}-U_n\|_{\TV}^2
 \le o(1)+C^{16}\sum_{\rho:\dim\rho>1}D^{2-16/4}
 =o(1)+C^{16}\sum_{\rho:\dim\rho>1}D^{-2}\longrightarrow0
\]
by Lemma~\ref{l:degree-inverse-square}. Replacing the eight independent
factors $\nu_n$ by $\mu_n$ costs at most
$8\|\nu_n-\mu_n\|_{\TV}=o(1)$.
\end{proof}

\subsection{A power expansion for 256 omitted blocks}

The second transfer uses the unnormalized average of good projections.
The preceding lemma supplies its input with $q=256$ and the same
threshold $D^{1/16}$. Here a lower operator bound on that average
controls the remainder, while expanding its complement costs the sum
of the absolute polynomial coefficients. This gives the stated, weaker
exponent without the scalar normalization used above.

For probabilities $p,q$ on a finite set, write
$D(p\Vert q)=\sum_xp(x)\log(p(x)/q(x))$ for relative entropy,
using natural logarithms, $0\log(0/q)=0$, and value $+\infty$
if $p$ gives positive mass where $q$ vanishes.

\begin{theorem}[Random-domain transfer]
\label{l:domain-transfer}
Let $n\to\infty$ through multiples of $256$, and let $\mu_n$ be
probabilities on $S_n$ satisfying
\[
 \delta_n=\E_{|R|=n/256}
       \|(\mu_n)_{-R}-U_{-R}\|_{\TV}\longrightarrow0.
\]
Condition on the event that at least $7/8$ of the sets $R$ satisfy
$\mu_n(gH_R)\le2/[S_n:H_R]$, and call the resulting law $\nu_n$.
For all sufficiently large $n$ it is defined, satisfies
$\|\nu_n-\mu_n\|_{\TV}\le16\delta_n$, and obeys
\[
 \|\widehat\nu_n(\rho)\|_{\op}\le CD^{-1/8}
 \qquad(D=\dim\rho>1)
\]
with an absolute $C$. If $\E_{\mu_n}\sgn(g)=o(1)$, then
$\|\mu_n^{*32}-U_n\|_{\TV}\to0$.
The averaged relative-entropy hypothesis
\[
 \E_{|R|=n/256}D((\mu_n)_{-R}\|U_{-R})=o(1)
\]
is sufficient for these conclusions.
\end{theorem}

\begin{proof}
Use Lemma~\ref{l:adaptive-domain-preparation} with $q=256$ and
$\xi=1/16$. In particular $c\ge15/16$. For a retained permutation put
\[
 M_g=\E_R e_R(g)P_R.
\]
Removing at most $1/8$ of the projections from their scalar average
gives
\begin{equation}
 \tfrac12 I\le\tfrac{13}{16}I\le M_g\le I.
 \label{l:domain-positive-average}
\end{equation}
For fixed unit vectors $u,v$ and every integer $p\ge1$,
\begin{equation}
 \E_\nu|\langle u,\rho(g)M_g^p v\rangle|
 \le AD^{-3/8},\qquad A=2\sqrt{C_2}.
 \label{l:domain-power-coefficient}
\end{equation}
Indeed, expand $M_g^p$ over $R_1,\ldots,R_p$ in that order.
After taking absolute values, retain only the indicator $e_{R_1}(g)$.
For the fixed sets, apply \eqref{l:adaptive-good-coefficient} to the
deterministic vector $z=P_{R_2}\cdots P_{R_p}v$, of norm at most one.
This proves \eqref{l:domain-power-coefficient}, even when the omitted
sets overlap.

The estimates for positive powers must now recover the original vector.
Put $L=\lfloor(\log_2D)/8\rfloor$. The identity
\[
 v=\sum_{j=0}^{L-1}M_g(I-M_g)^jv+(I-M_g)^Lv
\]
holds also when $L=0$, with an empty sum. Its remainder has norm at
most $2^{-L}$ by \eqref{l:domain-positive-average}. Expand each
$M_g(I-M_g)^j$ as a polynomial in $M_g$. The absolute values of its
coefficients sum to $2^j$, so \eqref{l:domain-power-coefficient} gives
\[
 \E_\nu|\langle u,\rho(g)v\rangle|
 \le2^{-L}+A\sum_{j=0}^{L-1}2^jD^{-3/8}
 \le2^{-L}+A2^LD^{-3/8}
 \le CD^{-1/8}.
\]
Here $2^{-L}\le2D^{-1/8}$ and $2^L\le D^{1/8}$, including $L=0$.
This proves the asserted operator bound.

The sign coefficient is again $o(1)$. For $32$ factors the remaining
Plancherel sum is at most
\[
 C^{64}\sum_{\rho:\dim\rho>1}D^{2-64/8}
 =C^{64}\sum_{\rho:\dim\rho>1}D^{-6}\longrightarrow0,
\]
because $D^{-6}\le D^{-2}$ and Lemma~\ref{l:degree-inverse-square}
applies. The total-variation cost of replacing the $32$ factors is
at most $32\|\nu_n-\mu_n\|_{\TV}=o(1)$.
Finally Pinsker's inequality and Jensen's inequality give
\[
 \delta_n\le
 \left(\tfrac12\E_R D((\mu_n)_{-R}\|U_{-R})\right)^{1/2},
\]
which proves the entropy version. In this normalization Pinsker's
inequality reads $\|p-u\|_{\TV}^2\le D(p\|u)/2$. To check the
factor, take $A=\{p\ge u\}$, set $a=p(A)$ and $b=u(A)$, and group
the entropy sum over $A$ and its complement. The log-sum inequality
gives $D(p\|u)\ge D(\operatorname{Ber}(a)\|
\operatorname{Ber}(b))$. As a function of $a$, the binary relative
entropy on the right vanishes with zero derivative at $a=b$ and has second
derivative $1/[a(1-a)]\ge4$. It is therefore at least
$2(a-b)^2=2\|p-u\|_{\TV}^2$, with boundary cases obtained by limits.
\end{proof}
